\begin{summarybox}
\textbf{\textcolor{CSummary}{10 秒速记：两端 $\pm a$，平方差出 $b^2$}}

椭圆 $\frac{x^2}{a^2}+\frac{y^2}{b^2}=1$（$a>b>0$），长轴端点切线 $x=\pm a$；第三条切线的切点 $P$ 不是 $(\pm a,0)$。

\[
\boxed{h_1h_2=b^2},\qquad
\boxed{S=a(h_1+h_2)\ge2ab},\qquad
\boxed{\text{取等 }P=(0,\pm b)}.
\]

\medskip
\textbf{\textcolor{CSummary}{最短现场推导}}

\[
\underbrace{y=kx+m}_{\text{动切线}},\quad
\underbrace{m^2=a^2k^2+b^2}_{\text{相切条件}},\quad
\underbrace{y_1=m-ak,\ y_2=m+ak}_{x=\mp a}.
\]
\[
\boxed{y_1y_2=m^2-a^2k^2=b^2},\qquad
\boxed{S=2a|m|\ge2ab}.
\]

\medskip
\textbf{\textcolor{CSummary}{忘记切线判定式时}}

设切点 $P(x_0,y_0)$，从 $\frac{x_0x}{a^2}+\frac{y_0y}{b^2}=1$ 出发；令 $x=\pm a$ 得
\[
h_1=\frac{b^2}{|y_0|}\left(1+\frac{x_0}{a}\right),\quad
h_2=\frac{b^2}{|y_0|}\left(1-\frac{x_0}{a}\right).
\]
\[
h_1h_2=b^2,\qquad
S=\frac{2ab^2}{|y_0|}\ge2ab.
\]
一句话：\textbf{“$1\pm x_0/a$：相乘平方差，相加得 $2$。”}

\medskip
\textbf{\textcolor{CSummary}{更直观的记忆}}

令 $P=(a\cos\theta,b\sin\theta)$，则两交点的纵坐标为
\[
y_1=b\cot\frac\theta2,\qquad
y_2=b\tan\frac\theta2.
\]
\textbf{$\tan$ 与 $\cot$ 互倒，所以乘积为 $b^2$；定积下等高最省面积。}

\medskip
\textbf{\textcolor{CSummary}{高考解题检查卡}}
\begin{itemize}[leftmargin=*,itemsep=0.3em,topsep=0.2em]
\item 识别：\textbf{端点处两条平行切线 + 第三条动切线}。
\item 图形：两条竖直线段是\textbf{梯形的底}，底间距 $2a$。
\item 细节：$h_i=|y_i|$；上下半椭圆都适用；\textbf{长轴端点排除}。
\item 最值：$h_1h_2=b^2$，基本不等式得到 $h_1+h_2\ge2b$。
\item 对称迁移：两条短轴端点切线 $y=\pm b$ $\Rightarrow$ \textbf{水平距离积 $a^2$}；面积仍最小 $2ab$。
\item 边界：切点趋向长轴端点时，面积趋向无穷，\textbf{不存在有限最大值}。
\end{itemize}
\end{summarybox}
